\chapter{Источники данных и контроль качества}

\section{Общие замечания}

Практический анализ данных в электроэнергетике начинается не с выбора модели, а с понимания того, откуда получены данные, в каком формате они хранятся и насколько им можно доверять.
Даже при корректной математической постановке задачи результат будет ненадежным, если исходные данные содержат систематические пропуски, дубликаты, смешение единиц измерения или неявные ошибки в формате.

\section{Типовые источники данных}

В инженерной практике наиболее часто используются следующие типы источников:
\begin{itemize}[leftmargin=1.25cm]
\item экспорт измерений из SCADA и АСУ ТП;
\item табличные выгрузки в форматах \code{CSV}, \code{XLSX}, \code{TXT};
\item журналы событий и эксплуатационные логи;
\item данные мониторинга режима и качества электроэнергии;
\item результаты лабораторных измерений и диагностических обследований;
\item выгрузки из корпоративных баз данных.
\end{itemize}

Для каждого набора данных необходимо фиксировать:
\begin{enumerate}[leftmargin=1.25cm]
\item источник происхождения;
\item формат файла;
\item состав признаков;
\item возможные ограничения и риски качества;
\item способ дальнейшего использования в анализе или моделировании.
\end{enumerate}

\section{Работа с разными форматами файлов}

В реальной инженерной практике данные редко поступают в одном унифицированном виде.
Даже в пределах одного исследовательского контура специалист может получить \code{CSV}-выгрузку измерений, файл \code{XLSX} с расчетными таблицами, текстовый файл \code{TXT} с разделителями и служебный \code{JSON}, описывающий состав файлов или параметры эксперимента.

Поэтому обучающийся должен уметь не только читать один тип таблицы, но и быстро определить, какой инструмент загрузки нужен в каждом конкретном случае.
В учебном репозитории для этой цели подготовлены локальные примеры:
\begin{itemize}[leftmargin=1.25cm]
\item \code{data/sample/formats/grid\_stability\_fragment.csv};
\item \code{data/sample/formats/combined\_cycle\_power\_plant\_fragment.xlsx};
\item \code{data/sample/formats/household\_power\_fragment.txt};
\item \code{data/sample/formats/dataset\_catalog.json}.
\end{itemize}

На их основе в связанном блокноте демонстрируется, что чтение разных форматов требует разных функций и параметров, однако последующая логика анализа остается общей: нужно проверить структуру, типы, полноту и смысловое назначение признаков.

\screenshotplaceholder
  {fig:file-format-examples}
  {Чтение файлов различных форматов в учебном блокноте}
  {Фрагмент блокнота, где последовательно показана загрузка файлов \code{CSV}, \code{XLSX}, \code{TXT} и \code{JSON} и выведены первые строки каждого источника.}
  {Открыть блокнот \code{02\_data\_sources\_and\_quality.ipynb}, выполнить раздел о форматах файлов и сделать скриншот так, чтобы были видны названия файлов, функции загрузки и первые строки таблиц.}

\figureinstruction
  {Примеры чтения тематических файлов разных форматов в едином учебном репозитории.}
  {После рисунка следует пояснить, что различие между форматами влияет на способ загрузки и параметры чтения, но не меняет требований к дальнейшему контролю качества и интерпретации данных.}

\section{Первичный контроль качества}

На первом этапе необходимо установить:
\begin{itemize}[leftmargin=1.25cm]
\item сколько строк и столбцов содержит набор данных;
\item какие типы имеют признаки;
\item присутствуют ли пропуски;
\item имеются ли дублирующиеся записи;
\item нет ли явных противоречий в шкалах и единицах измерения.
\end{itemize}

Даже эти простые проверки позволяют обнаружить существенные проблемы.
Так, повторные записи искажают группировки и средние значения, а некорректно распознанные типы времени могут разрушить весь дальнейший временной анализ.

\screenshotplaceholder
  {fig:data-quality-profile}
  {Первичный профиль набора данных в Jupyter-блокноте}
  {Таблица в блокноте, где показаны типы признаков, доля пропусков и число уникальных значений.}
  {Открыть блокнот \code{02\_data\_sources\_and\_quality.ipynb}, выполнить ячейки профилирования и сделать скриншот итоговой таблицы так, чтобы были читаемы названия столбцов и доли пропусков.}

\figureinstruction
  {Первичный профиль набора данных: типы столбцов, доля пропусков и число уникальных значений.}
  {После рисунка необходимо пояснить, что даже такой компактный профиль уже позволяет понять, какие признаки пригодны для непосредственного анализа, а какие требуют преобразования или очистки.}

\section{Паспорт датасета}

Для учебной и инженерной практики полезно оформлять краткий паспорт набора данных.
В него рекомендуется включать:
\begin{itemize}[leftmargin=1.25cm]
\item название и источник;
\item формат хранения;
\item масштаб набора данных;
\item целевую переменную, если задача уже определена;
\item ключевые признаки;
\item ограничения и риски;
\item краткий план предобработки.
\end{itemize}

Такой паспорт удобен не только для учебных целей, но и как элемент документации реального анализа.
\section{Что должен сделать обучающийся}

В результате выполнения главы обучающийся должен уметь:
\begin{itemize}[leftmargin=1.25cm]
\item прочитать табличный файл в Jupyter-блокноте;
\item описать набор данных в форме краткого паспорта;
\item выделить признаки с пропусками и дубликаты;
\item сформулировать первичный план очистки и подготовки данных.
\end{itemize}
